\documentclass[12pt, a4paper]{article}
\usepackage[utf8]{inputenc}
\usepackage{graphicx}
\usepackage{geometry}
\usepackage{amsmath}
\usepackage{amssymb}
\usepackage{booktabs}
\usepackage{float}
\usepackage{hyperref}
\usepackage{authblk}

\geometry{margin=1in}

\title{\textbf{Adaptive Confidence-Aware and Trust-Based EEG Framework for Multilevel Stress Detection, Uncertainty Estimation, and Selective Prediction: A Comprehensive Literature Review}}

\author{
    \textbf{Mr. M. Pruthvi Raj$^{1, a)*}$, T. Ranjithkumar$^{2, b)}$, G. Harika$^{3, c)}$, D. PriyaChandana$^{4, d)}$, Ch. Babiji$^{5, e)}$} \\
    \vspace{0.2cm}
    $^{1}$Sr. Asst. Prof, Dept. of CSE (AI\&ML), Lakireddy Bali Reddy College of Engg., Mylavaram, A.P, India \\
    $^{2,3,4,5}$Student, Dept. of CSE (AI\&ML), Lakireddy Bali Reddy College of Engg., Mylavaram, A.P, India \\
    \vspace{0.2cm}
    \small $^{a)}$jntupruthvi@gmail.com, $^{b)}$ranjithkumar3742@gmail.com, $^{c)}$harikagajj12@gmail.com \\
    \small $^{d)}$chandudevireddy67@gmail.com, $^{e)}$babjichode7@gmail.com
}
\date{}

\begin{document}

\maketitle

\begin{abstract}
Cognitive stress and mental workload detection via Electroencephalogram (EEG) signals play an essential role in neuroergonomics, human-computer interaction, and clinical diagnostics. Timely diagnosis of abnormal cognitive strain prevents mental fatigue and operational errors. Traditionally, extracting features from non-stationary EEG signals was heavily reliant on manual feature engineering. Recent progress in deep learning has introduced highly efficient automated methods for diagnosing mental states directly from raw or minimally processed EEG sequences. This massive and comprehensive literature review highlights the deep learning models developed for the detection and severity assessment of cognitive stress. Architectures such as Convolutional Neural Networks (CNNs), Long Short-Term Memory (LSTM) networks, and their hybrid variants are extensively analysed. Furthermore, the review explores widely used EEG datasets, performance evaluation metrics, and critical challenges including signal artefacts, subject-to-subject variability, and the lack of model interpretability and trust.
\end{abstract}

\vspace{0.5cm}
\noindent\textbf{Keywords:} Electroencephalogram (EEG), Deep Learning, Cognitive Stress, Convolutional Neural Networks, Trust-Based AI, Mental Workload, Uncertainty Estimation, Selective Prediction.

\newpage
\tableofcontents
\newpage

\section{Introduction}
Mental stress is a physiological and psychological response to high cognitive workload, often leading to fatigue, reduced decision-making capabilities, and long-term health complications. EEG is widely used for cognitive stress monitoring as it provides high temporal resolution recordings of brain electrical activity. As illustrated in various studies, shifts in specific frequency bands (such as Alpha suppression and Beta enhancement) are powerful indicators of mental exertion. However, examining multi-channel EEG recordings manually is extremely labour-intensive and prone to subjective variations. This has promoted the introduction of computer-aided deep learning systems capable of automated real-time mental state classification. The deployment of reliable, confidence-aware models is necessary to bridge the gap between academic prototypes and clinical reality.

\section{Methodology Pipeline}
\textit{(Note: Insert Flowchart Diagram Here mirroring the structure: Raw EEG $\rightarrow$ Artefact Removal $\rightarrow$ Feature Extraction $\rightarrow$ Deep Learning Model $\rightarrow$ Cognitive State Prediction)}

The development of Artificial Intelligence (AI) and, in particular, of its branches like Machine Learning (ML) and Deep Learning (DL) has made a significant contribution to automated biosignal analysis. Deep learning models enable the learning of discriminative temporal and spatial features without the need for hand-crafted feature extraction techniques, thus advancing the state-of-the-art in EEG analysis.

\begin{figure}[H]
    \centering
    \includegraphics[width=0.8\textwidth]{fig1.jpg}
    \caption{Flow chart of Cognitive Stress Detection Pipeline.}
\end{figure}

\section{Deep Learning Techniques for EEG Analysis}
Different deep neural network architectures like CNN, RNN, and hybrid combinations have been used for cognitive workload assessment. Table 1 summarises the deep learning models used for EEG analysis, highlighting their architectures, main advantages, and limitations.

\begin{table}[H]
\centering
\caption{Overview of Deep Learning Architectures for EEG Detection}
\vspace{0.2cm}
\resizebox{\textwidth}{!}{
\begin{tabular}{p{2cm} p{4cm} p{4cm} p{4cm}}
\toprule
\textbf{Model} & \textbf{Architecture Feature} & \textbf{Main Advantage} & \textbf{Limitation} \\
\midrule
\textbf{CNN} & Convolutional layers learn spatial features across channels & Automatic spatial feature extraction & Struggles with long-term temporal dependencies \\
\textbf{RNN / LSTM} & Recurrent gates and memory cells & Captures sequential and temporal dynamics & Computationally intensive, vanishing gradients \\
\textbf{Hybrid (CNN-LSTM)} & Spatial convolution followed by temporal recurrence & Strong spatio-temporal representation & High computational and memory requirements \\
\textbf{Attention} & Self-attention mechanism & Captures global temporal dependencies & Requires substantial training data \\
\bottomrule
\end{tabular}
}
\end{table}

\subsection{Convolutional Neural Networks (CNNs)}
CNNs are among the most widely used techniques for EEG classification. By treating the 2D EEG matrices (channels $\times$ time) as images, CNNs extract hierarchical features, effectively capturing the spatial correlations between different electrodes. 1D CNNs, in particular, are exceptional at convolving across the temporal dimension independently for each channel.

\subsection{Long Short-Term Memory Networks (LSTMs)}
Since EEG signals are continuous time-series data, LSTMs are particularly well-suited for analysing them. LSTMs utilise memory cells to retain information over long sequences, solving the vanishing gradient problem found in standard RNNs. The mathematical foundation of LSTMs includes the input gate, forget gate, cell state, and output gate, allowing the network to explicitly control the flow of temporal memory.

\subsection{Hybrid Architectures}
Recent literature demonstrates that combining CNNs for spatial feature extraction with LSTMs for temporal sequence modelling yields superior results in detecting subtle cognitive shifts in continuous EEG streams. By appending a self-attention mechanism, the network can weigh specific temporal epochs more heavily than others, creating an Adaptive Confidence-Aware representation.

\section{Datasets}
Available EEG datasets are critical for benchmarking models. Frequently used datasets include:
\begin{itemize}
    \item \textbf{STEW (Simultaneous Task EEG Workload):} Contains 14-channel EEG recordings from 48 subjects across resting and high cognitive workload states (the SIMKAP test). This is one of the most widely used datasets for evaluating mental workload classification models.
    \item \textbf{SEED:} Focuses on emotion and stress detection across multiple sessions.
    \item \textbf{DEAP:} A multimodal dataset for the analysis of human affective states, heavily utilized in literature for baseline comparisons.
\end{itemize}

\section{Evaluation Metrics}
The performance of deep learning models for EEG classification is analysed via accuracy, precision, recall, and F1-score. 

\begin{table}[H]
\centering
\caption{Evaluation Metrics Used for Assessing Deep Learning Models}
\vspace{0.2cm}
\begin{tabular}{p{3cm} p{6cm} p{5cm}}
\toprule
\textbf{Metric} & \textbf{Formula} & \textbf{Purpose} \\
\midrule
\textbf{Accuracy} & (TP+TN) / (TP+TN+FP+FN) & Overall proportion of correctly classified windows \\
\textbf{Precision} & TP / (TP+FP) & Measures false positive rate \\
\textbf{Recall} & TP / (TP+FN) & Shows how well the model spots true stress cases \\
\textbf{Selective Accuracy} & Accuracy calculated only on confident predictions & Evaluates the trust and reliability of the model \\
\bottomrule
\end{tabular}
\end{table}

The concept of Selective Accuracy, heavily emphasized in recent confidence-aware literature, operates by dividing the dataset into `accepted` and `abstained` windows. The metric ensures that high-risk, noisy windows are rejected, drastically boosting the operational safety of the system.

\section{Challenges in Deep Learning-Based EEG Detection}
Despite significant achievements, the actual deployment of these models encounters numerous challenges.

\subsection{Signal Noise and Artefacts}
EEG signals have a very low signal-to-noise ratio. They are easily contaminated by ocular artefacts (blinking), electromyogenic noise (muscle movement), and external electrical interference. 

\subsection{Subject Variability}
Models trained on one set of subjects often generalise poorly to new, unseen subjects due to immense inter-subject variability in skull thickness, brain folding, and cognitive baseline states. Transfer learning and subject-invariant feature learning are heavily researched to overcome this issue.

\subsection{Overconfidence and Trust}
Deep neural networks are often poorly calibrated, displaying high confidence even when predicting incorrectly on noisy data. Establishing Trust-Aware systems that can quantify uncertainty (e.g., via Monte Carlo Dropout) and abstain from predicting is a critical ongoing challenge.

\section{Conclusion}
Cognitive stress detection via EEG is a vital component of modern neuroergonomics. The current literature review analysed the use of deep learning algorithms for automated diagnosis. While hybrid CNN-LSTM networks have shown immense potential, future research must address inter-subject variability and focus on developing confidence-aware, transparent models suitable for real-world deployment. The introduction of mechanisms like Temperature Scaling and Monte Carlo Dropout represent a massive paradigm shift towards trustworthy Artificial Intelligence.

\end{document}
